\title{Large Language Models Can Automatically Engineer Features for Few-Shot Tabular Learning}

\begin{abstract}
Large Language Models (LLMs) have demonstrated impressive proficiency in addressing complex and unfamiliar reasoning tasks, signaling substantial promise for their application in tabular learning—a domain crucial to various practical problems. This study introduces \textsf{FeatLLM}, a novel in-context learning framework that leverages LLMs as feature engineers to craft a set of input features tailored for superior tabular prediction tasks. The features derived from this approach are subsequently utilized by a straightforward downstream machine learning model, such as linear regression, resulting in robust performance even in scenarios with limited labeled data. Notably, at prediction time, \textsf{FeatLLM} only relies on this simple predictive model and the identified features, dispensing with the need for repeated LLM access per inference instance. Relative to existing LLM-driven methodologies, \textsf{FeatLLM} avoids real-time querying of the LLM during inference and requires only API-level interaction with LLMs, effectively sidestepping restrictions related to prompt length. Empirical results over a broad spectrum of tabular datasets spanning multiple domains illustrate that \textsf{FeatLLM} creates high-caliber rules, consistently surpassing other approaches like TabLLM and STUNT by a significant margin (averaging a 10\% improvement).
\end{abstract}

\section{Introduction}
Large Language Models (LLMs) have recently exhibited exceptional capability in formulating effective responses to previously unencountered tasks, all without the necessity for explicit task-dependent fine-tuning~\cite{creswell2022selection,imani2023mathprompter,taylor2022galactica}. Owing to training on expansive textual datasets, LLMs capture extensive prior knowledge, enabling them to often replace human expert insight in various fields~\cite{Genomes2021,singhal2023large,wilcox2003role}, thus becoming integral in automating a wide range of applications, such as conversational analysis~\cite{finch2023leveraging} and automated program repair~\cite{fan2023automated}. LLMs are especially powerful when furnished with both task descriptions and a handful of representative samples within their prompts, allowing them to apprehend task context and focus selectively on relevant features and examples~\cite{brown2020language,zhao2023survey}. This underlines LLMs’ proficiency in data analysis and establishes their suitability for automated feature engineering.

Efforts to harness LLMs' embedded knowledge and reasoning have recently extended to tabular learning. For example, domain-specific information—such as the heightened susceptibility of older adults to certain diseases—can bolster predictive modeling for health outcomes. Contemporary strategies in this field involve expressing tasks and feature semantics in natural language, organizing the tabular information sequentially, and presenting it to an LLM~\cite{dinh2022lift,wang2023anypredict}. Subsequent procedures might include the use of in-context learning~\cite{nam2023semi} or the application of parameter-efficient fine-tuning~\cite{dinh2022lift,hegselmann2023tabllm}. By leveraging the domain knowledge encoded within LLMs, these methods have demonstrated improved performance relative to traditional tabular learning, particularly in data-scarce scenarios~\cite{hegselmann2023tabllm}.

Existing tabular learning approaches that rely on LLMs are subject to several important drawbacks. End-to-end prediction with these models demands at least one inference per individual data instance, significantly increasing computational load. Achieving strong predictive performance typically hinges on fine-tuning the LLM, which presents substantial practical barriers for models whose training resources or tuning interfaces are inaccessible; this issue is increasingly pronounced as high-performing recent LLMs often only offer restricted inference API access~\cite{anil2023palm,openai2023gpt4}. Additionally, the effectiveness of many of these methods diminishes with longer prompts~\cite{liu2023lost}; as the number of features in the tabular dataset expands, resulting in longer serialized prompts, such techniques either become unworkable or demonstrate a marked decrease in performance. Collectively, these limitations substantially restrict the scalability and practicality of LLM-driven tabular learning in real-world deployments.

To address these issues, our methodology reconceptualizes the function of LLMs in tabular learning. Rather than using LLMs solely for direct end-to-end inference, we employ them for advanced feature engineering, capitalizing on their pre-existing knowledge base more efficiently. We introduce a new in-context learning framework, \textsf{FeatLLM}, designed to uncover the “criteria” that inform LLM-based predictions. For example, in the context of disease prediction, the LLM is capable of articulating explicit decision rules that link specific combinations of feature values to the diagnosis. These generated criteria serve to construct novel features that supplant the original ones in predictive modeling tasks. With this paradigm, LLMs primarily facilitate feature generation, simplifying the downstream model—once new features are synthesized, inference can be executed by a lightweight model, thereby reducing inference time. Utilizing the in-context learning prowess of LLMs, feature extraction can occur without further training—demonstrative examples sufficing within the input prompt. Moreover, adopting ensemble-inspired strategies~\cite{breiman2001random}, such as feature bagging, allows us to alleviate issues arising from excessive prompt lengths.

\textsf{FeatLLM} decomposes the process of engineering novel features from prompts into two sequential sub-tasks: (1) comprehending the given problem and deducing the associations between features and targets; and (2) employing insights gathered from the initial stage, along with a set of few-shot examples, to derive critical decision rules that enable class separation. This disciplined, stepwise approach assists LLMs in disregarding irrelevant features during the rule construction phase. The formulated rules are subsequently executed on the dataset (interpreting them as programmatic code), leading to a binary transformation indicating for each rule whether a sample conforms. Afterwards, a simple model is trained to estimate class probabilities based on these synthesized features. To further enhance stability, an ensemble technique is utilized, generating features repeatedly in combination with feature bagging. This also alleviates the issue of prompt size limitations by calibrating the number of features and samples sampled during bagging. Free from the need for additional training and offering efficient inference, \textsf{FeatLLM} enables practical LLM applications across diverse real-world domains.

We conduct an extensive assessment of \textsf{FeatLLM} across 13 diverse tabular datasets under low-data conditions, demonstrating its consistent and competitive performance. Our experiments reveal that LLMs are adept at leveraging both pre-existing knowledge and dataset-specific information to construct high-value features. The proposed framework achieves superior results compared to state-of-the-art few-shot learning baselines in various experimental scenarios. The implementation is publicly accessible via anonymized GitHub at~\texttt{https://github.com/Sungwon-Han/FeatLLM}.

\section{Related Work}

\subsection{Few-Shot Learning with Tabular Data} Few-shot learning has made it possible to derive broadly applicable representations from data, despite having only limited annotated examples, thereby significantly reducing the labeling burden~\cite{chen2018closer}. Although the field began with a primary emphasis on image-based tasks, recent studies have established its adaptability and effectiveness in other domains, such as text, audio, and, increasingly, tabular data~\cite{majumder2022few,nam2023stunt,schick2021exploiting}. Learning from sparse supervision raises concerns about picking up coincidental or non-generalizable patterns~\cite{bartlett2021deep}. To counteract this, contemporary approaches integrate auxiliary data or leverage prior knowledge that is specific to the domain in question, thus embedding appropriate inductive biases into the model. Strategies include harnessing ample unlabeled data sets with carefully chosen augmentation procedures within a semi-supervised learning context~\cite{ucar2021subtab,yoon2020vime}, or securing robust initial model states via unsupervised, task-agnostic pre-training~\cite{somepalli2021saint}. Unsupervised learning often uses objectives such as reconstructing masked or perturbed sections of the data~\cite{arik2021tabnet,majmundar2022met}, or applying data augmentations to facilitate contrastive learning~\cite{bahri2022scarf,chen2020simple}. Nevertheless, the inherently mixed-type and diverse characteristics of tabular datasets complicate the discovery of data augmentations that are widely effective, resulting in limited gains in the few-shot regime~\cite{nam2023stunt}. 

To overcome these challenges, recent research has explored the potential of pretraining models on a large and diverse array of tabular datasets, far beyond those directly related to the end task, followed by transfer learning to the specific target problem~\cite{zhang2023generative,levin2022transfer,zhu2023xtab}. TransTab~\cite{wang2022transtab}, for example, leverages transformer models that not only process cell values but also capture semantic relationships by incorporating column names. The generalized knowledge acquired in this way supports strong zero-shot and few-shot performance when models are presented with novel tabular problems. Another example is TabPFN~\cite{hollmann2022tabpfn}, which utilizes mixtures of distribution types to generate synthetic training data for Bayesian neural network pretraining. Inference is performed by jointly presenting labeled and unlabeled data to the network without further optimization, allowing the transfer of learned inductive biases to new tasks. These methods have demonstrated superior results compared to traditional ensemble tree-based models, with particular advantage observed in few-shot scenarios.

\subsection{Language-Interfaced Tabular Learning} An emerging avenue for incorporating prior knowledge into tabular modeling involves harnessing large language models (LLMs)~\cite{anil2023palm,openai2023gpt4}. LLMs, owing to their training on vast and heterogeneous text datasets, are capable of strong generalization to novel tasks and have been successfully applied across a spectrum of applications~\cite{brown2020language}. There is growing interest in leveraging the latent knowledge within LLMs for tabular data applications. Typical approaches involve transforming tabular inputs into textual representations and integrating them into LLM prompts~\cite{dinh2022lift,hegselmann2023tabllm,wang2023anypredict}. By providing the models with not only the tabular data itself but also metadata such as feature and task descriptions within prompts, models can reason over both features and task requirements to make predictions. Advancements in this area include tailored model tuning with efficient adaptation techniques such as LoRA~\cite{hu2021lora} or IA3~\cite{liu2022few}, as well as prompt-based, training-free in-context learning by supplying a handful of exemplar cases within the prompt~\cite{dinh2022lift,hegselmann2023tabllm,nam2023semi,slack2023tablet}.

While previous approaches have proven valuable for enhancing few-shot tabular learning, their reliance on consistently routing inference through the LLM introduces barriers for use in industrial contexts. This work seeks to shift away from the typical end-to-end, black-box paradigm that invokes the LLM for every individual sample’s prediction. Rather, the focus is on leveraging the LLM to deduce explicit rules or conditions that define each class. \textsf{FeatLLM} translates these generated rules into programmatic code to derive new features, thus allowing subsequent predictions to bypass the LLM entirely. Through in-context learning, the model utilizes both prior knowledge and the limited available samples to distill actionable rules.

\section{Methods}
\textsf{FeatLLM} is constructed to identify and extract ``rules"—the attribute criteria tied to each answer class—by drawing from the limited provided examples, as visualized in Figure~\ref{fig:main_model}. These rules are interpreted by an LLM and subsequently transformed into binary indicators for each input, marking whether a given rule is met. The ensemble of these binary indicators forms a new feature vector, which is combined using non-negative trainable weights through a dot product operation to calculate each class's probability. Training this architecture enables the model to assess the contribution of every feature towards class membership, as further described in Section~\ref{Sec:training}. This methodology is iteratively applied with bagging, and final outputs are aggregated through ensembling. A detailed description of the problem formulation and each phase of the pipeline is provided below.

\vspace*{-0.12in}
\paragraph{Problem formulation.} Consider a tabular dataset $\mathcal{D} = \{(\mathbf{x}^i, \mathbf{y}^i)\}_{i=1}^{N}$ comprised of $N$ labeled instances. Each example $\mathbf{x}^i$ represents a $d$-dimensional vector, corresponding to $d$ input features, where $F = \{f_j\}_{j=1}^{d}$ lists the natural language names for these features, with definitions available independently. In a classification scenario, each label $\mathbf{y}^i$ is drawn from a finite collection of classes. For $k$-shot learning experiments, a subset of $k$ ($<N$) randomly chosen labeled instances is used for model training.

\subsection{Prompt Design for Extracting Rules}
\label{Sec:prompt_design}
To enable rule extraction through large language models (LLMs) along interpretable lines of reasoning, we structure the prompt to emulate the systematic thinking an expert might apply when tackling tabular data problems. The prompt integrates three principal elements (see Fig.~\ref{fig:prompt_for_rule} and Appendix~\ref{sec:example_rule}):

\vspace*{-0.12in}
\paragraph{Basic information description.} This segment presents the core context necessary for task resolution (highlighted in orange in Fig.~\ref{fig:prompt_for_rule}). It comprises: (1) a clearly stated task, together with label options; (2) concise feature descriptions, including their value types and, if available, their definitions; and (3) illustrative example demonstrations. Task descriptions take the form of explicit questions (e.g., ``Does this patient's myocardial infarction complications data indicate chronic heart failure? Yes or no?"; additional examples appear in Appendix~\ref{sec:task_desc}). Feature descriptions specify variable types and, for categorical attributes, may include representative category examples. Demonstrative instances serialize selected training data along with their true labels in a text-based format, adopting serialization conventions found in earlier works~\cite{dinh2022lift,hegselmann2023tabllm}:

\begin{align}
&\text{Serialize}(\mathbf{x}^i,\mathbf{y}^i, F) = \nonumber \\ 
&``\ f_1\ \text{is}\ \mathbf{x}^i_1.\ ... \ f_d\  \text{is}\  \mathbf{x}^i_d.\ \ \text{Answer:}\  \mathbf{y}^i\ ”
\end{align}

\vspace*{-0.12in}
\paragraph{Reasoning instruction.} Our objective is to strengthen the LLM’s reasoning capabilities by incorporating explicit reasoning instructions (highlighted in blue in Fig.~\ref{fig:prompt_for_rule}). These instructions start with an introductory statement akin to the chain-of-thought framework~\cite{wei2022chain}. We further segment the LLM’s reasoning into two distinct phases. Initially, the LLM is prompted to determine the causal links or prevailing trends between features and the task description. This step is performed without leveraging any example demonstrations, instead utilizing general domain knowledge or common sense to maximize the model’s retrieval of relevant prior knowledge. In the subsequent phase, the LLM draws on both the examples provided and its earlier insights to formulate class-specific rules. This staged reasoning discourages the model from focusing on irrelevant columns and the potential discovery of misleading correlations, while also emphasizing more influential features. To strike a balance and avoid generating either an insufficient or excessive number of rules—which might lead to underfitting or an unwieldy prompt—we mandate a set quota of rules to be generated per class (i.e., 10)\footnote{Analysis on the number of rules can be found in Section~\ref{sec:analysis}.}. Furthermore, to reduce the risk of rule-type mismatches, we instruct the LLM to reference the feature types given in the basic information section while crafting rules.

\vspace*{-0.12in}
\paragraph{Response instruction.} For improved interpretability and usability of the generated rules, we explicitly instruct the LLM on how to structure its outputs (see yellow highlights in Fig.~\ref{fig:prompt_for_rule}). The prompt outlines the required answer format for each class (i.e., format for the response) and the template that individual rules should adhere to (i.e., format for the rule). This approach allows the LLM to tailor its outputs appropriately, taking feature type into account. Additionally, unless otherwise directed, the LLM is permitted to synthesize compound rules through the use of logical connectors such as AND and OR. Illustrative examples of such responses can be found in Appendix~\ref{sec:outcome-rule}.

\subsection{Inferring Class Likelihood via Rules}
\label{Sec:training}

\paragraph{Parsing rules for feature generation.} The rules identified during the preceding phase are leveraged to construct new binary features. For each class, these features indicate whether a given sample meets the conditions of the relevant rules, resulting in binary values ('0' or '1'). These binary indicators can serve as inputs to estimate the likelihood of class membership for each sample. However, as the rules produced by the LLM are expressed in natural language, there is a need for an effective parsing mechanism to automate feature generation. Although formatting and response structure are constrained during rule generation to facilitate easier parsing, the LLM output may still contain syntactic inconsistencies or noise~\cite{kovavc2023large}. For example, the LLM may rephrase comparative operators, such as replacing symbols like ``$>$” or ``$<$” with verbal descriptions like ``greater than" or ``less than." Such inconsistency complicates the parsing process.

Rather than devising intricate parsing code to handle noisy and variable natural language output, we employ the LLM itself to tackle this challenge. Owing to its proficiency in comprehending semantic content and generating executable code (thanks to exposure to code data), the LLM can convert these rules into parseable logic. To accomplish this, we design prompts that provide the function name, specifications for input and output, as well as the derived rules, and submit these prompts to the LLM. Illustrative prompts and the resulting outputs are provided in Figures~\ref{fig:prompt_for_parsing} and ~\ref{sec:example_parsing}-\ref{sec:example_parse_outcome} in the Appendix. The code output by the LLM is run via Python’s \texttt{exec()} method to facilitate the data transformation process.

\vspace*{-0.12in}
\paragraph{Inferring class likelihood.} With the introduction of these rule-derived binary features for each class, an intuitive approach for estimating the class likelihood is to compute, for each class, the number of corresponding rules satisfied by the sample—formally, summing the relevant binary features. Nonetheless, since rules can vary in their predictive value, it becomes important to learn the weight or importance of each rule from labeled training data. To this end, we fit a standard machine learning (ML) model using the binary features as input for each class separately, enabling the model to learn optimal weights. As an example, a bias-free linear model can be utilized. Let $\mathbf{z}_k^i \in \{0, 1\}^R$ denote the binary feature vector corresponding to sample $\mathbf{x}^i$ and class $k$, where $R$ is the total number of rules for that class, and $c$ refers to the total number of classes. The predicted probability vector $\mathbf{p}^i$ for each class is then computed as follows:

\begin{align}
\text{logit}_k^i &= \max(\mathbf{w}_k, 0) \cdot \mathbf{z}_k^i, \nonumber \\
\mathbf{p}^i &= \text{Softmax}({[\text{logit}_{1}^i, ..., \text{logit}_{c}^i]}). \label{eq:infer_prob}
\end{align}

In this context, $\mathbf{w}_k$ denotes the adjustable parameters within the linear framework, reflecting the relevance of each individual feature. Restricting these weights to a non-negative space guarantees that features produced by the LLM do not detract from the likelihood assigned to any class during training. The resulting logit vector is subsequently transformed into class probabilities through the softmax function. Optimization of $\mathbf{w}_k$ is achieved via cross-entropy loss, utilizing a limited set of training examples in a few-shot regime.

\vspace*{-0.12in}
\paragraph{Ensembling with bagging.} To improve predictive accuracy, the entire inference process is executed multiple times, resulting in an ensemble of models for final decision-making. Specifically, by assembling the set of learned weights from $T$ separate iterations, represented as $\{\mathbf{w}[t]\}_{t=1}^T$, the final prediction is obtained by averaging the per-class probabilities $\mathbf{p}^i$ derived from each individual weight vector $\mathbf{w}[t]$ as formalized in Eq.~\ref{eq:infer_prob}.

To foster diversity in the rule sets generated for the ensemble, a deliberately elevated temperature parameter is chosen during LLM inference. Additionally, the sequence of few-shot demonstrations is shuffled, drawing on findings that this order can substantially affect LLM outputs~\cite{lu2022fantastically}\footnote{Further analysis of rule diversity is presented in Appendix~\ref{sec:diversity}}. To capitalize on the diversity arising from ensemble predictions and thus bolster robustness, bagging is utilized, where random subsets of features or data instances are selected for each trial. This becomes particularly beneficial when handling larger volumes of training data or high-dimensional feature spaces. The advocated ensemble method offers two principal benefits. First, any erroneous rule sets produced in certain trials may be counterbalanced by correct rules from alternative trials, thus fortifying the framework’s robustness. This leverages the LLM’s self-consistency property~\cite{wang2022self}, where rules repeatedly inferred across independent runs have a higher chance of correctness. Second, the ensemble technique mitigates the prompt-length limitation inherent to LLMs. In situations where tabular data surpass prompt length constraints, the use of bagging helps maintain manageable input sizes while sustaining robust model performance. Since any individual set of rules might not exhaustively capture feature relationships, aggregating predictions from several weak learners assembled over multiple runs offsets the shortcomings of incomplete feature or instance representation within singular trials.

\section{Experiments}
We assess the effectiveness of \textsf{FeatLLM} using a variety of tabular datasets, and carry out several analyses to explore the mechanisms underlying our framework. Further experimental results, such as those involving alternate LLMs, qualitative assessments, and other details, are provided in the Appendix due to limitations of space.

\subsection{Performance Evaluation}
\paragraph{Datasets.}
Our study is conducted on 13 datasets encompassing binary and multi-class classification tasks: (1) Adult~\cite{asuncion2007uci}, which involves predicting if a person's yearly income surpasses \$50,000; (2) Bank~\cite{moro2014data}, focusing on the likelihood that a client subscribes to a term deposit; (3) Blood~\cite{yeh2009knowledge}, concerning prediction of whether blood donors will return; (4) Car~\cite{kadra2021well}, which categorizes the quality rating of automobiles; (5) Communities~\cite{misc_communities_and_crime_183}, designed to predict regional crime rates; (6) Credit-g~\cite{kadra2021well}, which assesses if a person represents a good or bad credit risk; (7) Diabetes\footnote{\texttt{kaggle.com/datasets/uciml/pima-indians-diabetes-database}}, for determining diabetes status in patients; (8) Heart\footnote{\texttt{kaggle.com/datasets/fedesoriano/heart-failure-prediction}}, focusing on forecasting the presence of coronary artery disease; and (9) Myocardial~\cite{misc_myocardial_infarction_complications_579}, for diagnosing chronic heart failure in individuals.

Additionally, we incorporate more contemporary datasets, along with synthetic ones, which have not been commonly referenced and were absent from the LLMs’ pre-training corpus: (10) Cultivars, which predicts the grain yield levels of soybean cultivars\footnote{\texttt{archive.ics.uci.edu/dataset/913}}; (11) NHANES, which predicts age group categorization from personal records\footnote{\texttt{archive.ics.uci.edu/dataset/887}}; (12) Sequence-type, assigning sequences to one of four categories (arithmetic, geometric, Fibonacci, or Collatz); and (13) Solution-mix, which determines if the concentration percentage in a mixture of four solutions exceeds 0.5. The first two of these datasets were only recently added to the UCI Machine Learning Repository after September 2023, ensuring their absence from the pre-training data of the LLM API (gpt-3.5-turbo-0613) utilized in our approach. The last two datasets are synthetically generated, making it impossible for the LLM API to have previously encountered or memorized their patterns prior to evaluation.

The datasets differ considerably in terms of scale and complexity, with further specifications outlined in Table~\ref{Tab:basic-desc}.
All datasets provide explicit names and detailed attribute descriptions. Datasets such as `Communities' and `Myocardial' are characterized by a large attribute set—exceeding 100 features—which creates significant challenges given the restricted context windows available to LLMs.

\vspace*{-0.12in}
\paragraph{Baselines.}
Our experimental comparison includes nine baseline methods. The initial three represent classical models for learning from tabular data: (1) Logistic Regression (LogReg), (2) XGBoost~\cite{chen2016xgboost}, and (3) RandomForest~\cite{ho1995random}. We further consider two methods that utilize unlabeled data: (4) SCARF~\cite{bahri2022scarf} and (5) STUNT~\cite{nam2023stunt}, although it is important to note that acquiring extensive unlabeled datasets is often impractical in real-world scenarios. Among more recent approaches, (6) TabPFN~\cite{hollmann2022tabpfn} stands out for generating diverse synthetic datasets to pretrain the model. Our evaluation also includes three LLM-based baselines: (7) In-context learning (In-context)~\cite{wei2022emergent}, (8) TABLET~\cite{slack2023tablet}, and (9) TabLLM~\cite{hegselmann2023tabllm}. In the case of in-context learning, a small number of labeled samples are embedded into the prompt, eschewing further model tuning. TABLET extends this approach by enriching the prompt with information from external classifiers—specifically, rules and prototypes—to sharpen inference quality. TabLLM, meanwhile, applies parameter-efficient tuning strategies such as IA3~\cite{liu202

\vspace*{-0.12in}
\paragraph{Implementation details.}
\textsf{FeatLLM} adopts GPT-3.5 as its primary LLM architecture, though its framework is not limited to this specific model (see Appendix~\ref{sec:backbones} for experiments with alternative backbones). The inference settings for the LLM include a temperature of 0.5 and a top-p parameter left at the default API value of 1. For the ensemble approach and rule extraction procedure, we fix the ensemble count to 20 and the number of extraction rules to 10. Additional exploration of hyper-parameter choices can be found in Figure~\ref{fig:hyperparameter2} and Appendix~\ref{sec:hyper-parameter}. 

For training the linear model, the Adam optimizer is selected with a learning rate of 0.01, running for a total of 200 epochs. Model selection leverages $k$-fold cross-validation to determine the optimal epoch. In establishing baselines, in-context learning is implemented using GPT-3.5, while TabLLM experiments employ T0~\cite{sanh2022multitask}, adhering to the setup specified by prior work. Importantly, TabLLM only permits the use of LLMs with open checkpoint availability, precluding the inclusion of GPT-3.5. For conventional machine learning models such as LogReg, XGBoost, and RandomForest, parameter tuning is conducted by combining Grid-search and $k$-fold cross-validation. Here, $k$ is chosen as either 2 or 4, ensuring the training data for each fold contains at least one instance from every class. Further implementation specifics are provided in Appendix~\ref{sec:baselines}.

\vspace*{-0.12in}
\paragraph{Results.}
Tables~\ref{tab:main1} and \ref{tab:main2} display comparative performance metrics across various datasets. Each experiment is conducted three times with different random partitions of training and test data, and we report the mean and standard deviation of the resulting AUC scores. \textsf{FeatLLM} either leads or places second among all compared baselines across the evaluated datasets. Remarkably, with as few as 4 shots, our approach matches the performance of traditional tabular baselines that utilize 32 shots (refer to Fig.~\ref{fig:compares}). Although STUNT and TabPFN achieve notable gains on some individual datasets, their aggregate improvement over conventional machine learning techniques remains limited. Additionally, LLM-driven methods such as in-context learning, TABLET, and TabLLM are unable to process tabular data with extensive feature sets. By contrast, our proposed method, leveraging bagging and ensembling, remains robust and performs well on high-dimensional datasets. While our bagging and ensemble strategy can, in principle, be incorporated into other LLM-based models, doing so would effectively double the number of inference calls per sample, significantly increasing computational requirements and rendering such adaptations infeasible.

\subsection{Analysis \& Discussion}
\label{sec:analysis}
\paragraph{Ablation study.} To assess the impact of individual components on overall performance, we conduct ablation studies targeting the following: (1) \textsf{-Tuning}, which excludes the linear model’s weight tuning, using a simple sum of binary feature outputs for logit computation instead; (2) \textsf{-Ensemble}, in which the ensemble phase is omitted; (3) \textsf{-Description}, where feature descriptions are not provided; and (4) \textsf{-Reasoning}, where the Step-1 element of the reasoning instruction for rule generation is removed.

Table~\ref{tab:ablation} reports the mean AUC change due to each ablation across all datasets. Performance deteriorates in every scenario when any of the key components are excluded. Notably, as the number of shots increases, the advantages of weight tuning become more pronounced. This pattern suggests that the effectiveness of tuning improves with more abundant data, as rule importance can be estimated more accurately. Conversely, omitting feature descriptions or reasoning instructions has a more substantial negative impact when few shots are available, highlighting the value of LLM prior knowledge in data-scarce regimes. The ensemble method, as introduced, consistently boosts results in every configuration.

\vspace*{-0.12in}
\paragraph{Training \& inference time comparison.} We assess and compare the durations required for both training and inference across the various methods. All tests utilize the Adult dataset, employing a training subset with 16 examples per class. By default, experiments use a single A100 GPU, with the exception of TabLLM, which requires four A100 GPUs due to model parallelization. The details of these runtime comparisons are listed in Table~\ref{tab:cost}. Of particular note, \textsf{FeatLLM} achieves notably efficient inference times that are on par with those of standard machine learning approaches. By contrast, inference times for alternative LLM-based techniques, such as In-context learning, TABLET, and TabLLM, are substantially greater. For training time, \textsf{FeatLLM} is largely similar to other LLM-based models, entailing just 30 API requests—a quantity that remains cost-effective and is amenable to parallel execution.

\vspace*{-0.12in}
\paragraph{Quality of features generated by \textsf{FeatLLM}.}
To evaluate the effectiveness of the rule-based features produced by \textsf{FeatLLM}, we design a straightforward experiment focused on their informativeness for practical applications. Concretely, we compute the mean AUROC for the features generated and their association with the target labels, further quantifying the proportion of features whose AUROC exceeds 0.5 within each dataset. Table~\ref{tab:quality} below presents an overview of these findings. The results demonstrate that the majority of generated rules pertain to the target variable and deliver meaningful information for task resolution (as indicated in the “Before tuning” column). Furthermore, throughout the linear model fine-tuning stage, those rules that display minimal relevance (i.e., rules with learned weights under the threshold) are systematically excluded (see the “After tuning” column). The threshold was fixed at 0.1, informed by the distribution of learned weights.

\vspace*{-0.12in}
\paragraph{Effect of adding spurious correlations.} To evaluate the ability of different approaches to eliminate irrelevant spurious correlations in scenarios with limited labeled data, we performed a targeted experiment. We chose the Heart dataset as the primary task and incrementally appended random columns taken from the Adult dataset—columns that bear no relation to the actual classification problem of the Heart dataset. The subsequent variation in model performance was then tracked as the number of extraneous columns increased. For the sake of consistency and equitable comparison, all models operated with only 8 labeled samples and no access to unlabeled data, leading us to omit approaches such as SCARF and STUNT.

Figure~\ref{fig:spurious} illustrates how model accuracies shift in response to injected spurious correlations. Among the approaches examined, \textsf{FeatLLM} demonstrated superior resilience, exhibiting the smallest drop in performance attributable to these noisy features. This robustness can be attributed to the framework’s enforcement of explicit reasoning about the connection between features and the task during rule extraction, drawing upon prior knowledge and thereby avoiding reliance on irrelevant features. In practice, rules referencing the noisy, randomly-injected columns made up merely 1-2\% of all extracted rules. Nonetheless, we also note that \textsf{FeatLLM} is not wholly immune—if additional columns drawn from data that is semantically or contextually related to the Heart dataset are appended, the model may still absorb spurious associations and confuse the underlying causality (see Appendix~\ref{sec:spurious-appendix}).

\vspace*{-0.12in}
\paragraph{Hyper-parameter analysis}
\textsf{FeatLLM} features two primary hyper-parameters: the ensemble count and the number of rules per class extracted per LLM invocation. To assess how these factors influence performance, we performed a series of evaluations. The findings indicate that increasing the ensemble size yields performance enhancements; however, such improvements plateau once the ensemble number exceeds roughly 20 (refer to Fig.~\ref{fig:appendix-hyperparameter1} in the Appendix). Concerning the number of rules, our experiments did not reveal a statistically significant difference, but we concluded that extracting 10 rules represents the best compromise between keeping the prompt concise and maintaining strong performance (see Fig.~\ref{fig:hyperparameter2}). We hypothesize that while a higher rule count augments input dimensionality and the informational content, it may also introduce overfitting, especially in settings with limited data.

\section{Conclusion}
We introduce \textsf{FeatLLM}, a method that sets a new benchmark for LLM-driven few-shot tabular data modeling. Distinct from prior methods that rely on LLMs for individual sample predictions, \textsf{FeatLLM} leverages LLMs in feature engineering by generating rules for predictive modeling. By incorporating rule creation and employing bagging for ensembling, \textsf{FeatLLM} offers low inference overhead, operates solely with API-based inference (no model retraining needed), and adapts flexibly to various data feature dimensionalities. As a result, \textsf{FeatLLM} attains robust predictive accuracy and serves as a competitive and data-efficient solution for real-world tabular applications. 

\textsf{FeatLLM} was specifically built for low-shot learning contexts. Future directions will focus on broadening its utility for datasets with larger sample sizes, innovating new feature engineering mechanisms beyond rule-based formulations, and enhancing interpretability, thereby enabling application to a more extensive set of problem domains.

\section{Broader Impact}
\textsf{FeatLLM} has been developed to harness the pretrained knowledge present in LLMs for superior tabular learning outcomes, aiming to surpass standard supervised learning results, particularly under data constraints. By injecting relevant prior knowledge into the learning process, our work directly tackles the problem of overfitting arising from sparse labeled datasets, promoting data efficiency. The advantages of \textsf{FeatLLM} are particularly pertinent for users in real-world domains, including finance and healthcare, where labeled data acquisition can be prohibitively costly or otherwise impractical, and the value of domain knowledge encoded within LLMs is substantial (as these models have been exposed to related domain-specific textual corpora). Looking forward, a critical avenue for broadening \textsf{FeatLLM}’s societal impact will be the investigation of potential societal biases and misinformation present in LLMs’ knowledge bases, as such factors could unduly influence—perhaps even outweigh—empirical evidence from labeled data.

\end{document}